
\section{$\omega_X$ 的\trAdjunctionFormula{}与 $\omega_X=\mathscr{K}_X$}

（在本节——全书最后一节——中，尽管我们曾在 \S29.1.6 中就术语提出过异议，我们仍将把「\trDualizingSheaf」一词用于 \trSerreDuality{}的\trCanonical{}层，原因在于：本节的全部要点正是要表明，「\trCanonical{}层」这一短语的两种用法在\trSmooth\trProjective\trVariety{} $X$ 的情形下指的是同一个东西：在 \S29.1.5 的 $\mathrm{Hom}$ 的\trSerreDuality{}中出现的那个层 $\omega_X$，以及 $\mathscr{K}_X=\det\Omega_X$。）

\trDualizingSheaf{}在用\trEffectiveCartierDivisor{}作\trSlicing{}时表现良好，而这对得到\mbox{\trDesideratum{}~29.1.1}——亦即 $\mathscr{K}=\det\Omega$ 是\trSerre{}对偶化的这一奇迹（\S18.5.2）——将会很有用。为了说明这一点，我们将对满足 $\mathrm{Ext}$ 的\trSerreDuality{}的\trVariety{}之\trSubvariety，给出其\trDualizingSheaf{}的一个描述。

作为热身，我们将在 $\pi$ 为\trClosedEmbedding{}的特殊情形下，把 $\pi^!$ 的概念推广到 $\OO$-\trModule{}上。

\textbf{29.4.A.（练习）}\quad 设 $\pi\colon X\to Y$ 是概形之间的一个\trClosedEmbedding，而 $\mathscr{G}$ 是一个 $\OO_Y$-\trModule。请解释为什么 $\mathscr{H}om_Y(\pi_*\OO_X,\mathscr{G})$ 自然地带有 $\OO_X$-\trModule{}的结构。提示：若 $\mathscr{I}$ 是 $X$ 的\trIdealSheaf，请解释 $\mathscr{H}om_Y(\pi_*\OO_X,\mathscr{G})$（在某个\trOpenSubset{} $\mathfrak{U}\subset Y$ 上）如何被「在 $X$ 上消失的函数」（$\mathscr{I}(\mathfrak{U})$ 中的元素）所零化。

\textbf{29.4.1.}\ \emph{\trhDefinition{}。}\quad 于是，若 $\pi\colon X\to Y$ 是一个\trClosedEmbedding，我们就已经定义了一个映射 $\pi^!\colon\mathit{Mod}_{\OO_Y}\to\mathit{Mod}_{\OO_X}$，它是\mbox{练习~17.1.K(a)} 中那个映射的扩张。在这种 $\pi$ 为\trClosedEmbedding{}的情形下，记 $\pi^!\mathscr{G}:=\pi^{-1}\mathscr{H}om(\pi_*\OO_X,\mathscr{G})$ 是合理且正确的，也就是说，在定义中使用符号 $\pi^{-1}$。

\textbf{29.4.B.（练习）}\quad 我们仍处于\mbox{练习~29.4.A} 的情形中。

(a) 证明：$(\pi_*,\pi^!)$ 是 $\mathit{Mod}_{\OO_X}$ 与 $\mathit{Mod}_{\OO_Y}$ 之间的一个\trAdjointPair。

(b) 证明：$\pi^!$ 把内射的 $\OO_Y$-\trModule{}送到内射的 $\OO_X$-\trModule。（提示：$\pi_*$ 是\trExact{}的；用\mbox{练习~23.5.B}。）

我们现在把这一讨论从 $\mathscr{H}om$ 推广到 $\mathscr{E}xt$。设 $\pi\colon X\to Y$ 是一个\trClosedEmbedding。那么，对任意 $\OO_Y$-\trModule{} $\mathscr{F}$，$\mathscr{E}xt^i_Y(\pi_*\OO_X,\mathscr{F})$ 自然地带有 $\OO_X$-\trModule{}的结构。（理由：我们通过如下方式计算它：取 $\mathscr{F}$ 的一个由 $\OO_Y$-\trModule{}构成的内射分解，加以截断，然后施加 $\mathscr{H}om_Y(\pi_*\OO_X,\cdot)$。而对任意 $\OO_Y$-\trModule{} $\mathscr{G}$，$\mathscr{H}om_Y(\pi_*\OO_X,\mathscr{G})$ 就带有 $\OO_X$-\trModule{}的结构，也就是 $\pi^!\mathscr{G}$。）为了强调它作为 $\OO_X$-\trModule{}的结构，我们把它写作 $\pi^{-1}\mathscr{E}xt^i_Y(\pi_*\OO_X,\mathscr{F})$。

\textbf{29.4.C.（练习）}\quad 证明：若 $Y$（从而 $X$）是\trLocally\trNoetherian{}的，而 $\mathscr{G}$ 是 $Y$ 上的一个\trCoherentSheaf，则 $\pi^{-1}\mathscr{E}xt^i_Y(\pi_*\OO_X,\mathscr{G})$ 是 $X$ 上的一个\trCoherentSheaf。（提示：$\mathscr{E}xt^i_Y(\pi_*\OO_X,\mathscr{G})$ 是 $Y$ 上的\trCoherentSheaf，\mbox{练习~29.2.F(b)}。）

\textbf{29.4.2. 用环绕簇的对偶化层表示子簇的对偶化层。}\quad 在本节余下的部分中，$\pi\colon X\to Y$ 都将是纯维\trProjective{} $k$-\trScheme{}之间的一个\trClosedEmbedding，二者的\trDimension{}分别为 $n$ 与 $N$，其中 $Y$ 满足对 $\mathrm{Ext}$ 的函子性\trSerreDuality。令 $r=N-n$（即 $X$ 在 $Y$ 中的\trCodimension）。由\mbox{推论~29.3.12}，$X$ 满足对 $\mathrm{Hom}$ 的函子性\trSerreDuality。我们现在要用 $\omega_Y$ 来识别 $\omega_X$。

\textbf{29.4.3. \trhTheorem。}\ ——\ \emph{我们有 $\omega_X=\pi^{-1}\mathscr{E}xt^r_Y(\pi_*\OO_X,\omega_Y)$。}

在证明\mbox{定理~29.4.3} 之前，我们先来解释，那些「先前的」$\mathscr{E}xt^i_Y(\pi_*\OO_X,\omega_Y)$ 都发生了什么。

\textbf{29.4.4. \trhProposition。}\ ——\ \emph{设 $\pi\colon X\hookrightarrow Y$ 是纯维\trProjective{} $k$-\trScheme{}之间的一个\trClosedEmbedding，二者的\trDimension{}分别为 $n$ 与 $N$，且 $Y$ 满足对 $\mathrm{Ext}$ 的函子性\trSerreDuality。则对一切 $i<r:=N-n$，都有 $\mathscr{E}xt^i_Y(\pi_*\OO_X,\omega_Y)=0$。}

\emph{证明。}\quad 由于 $\mathscr{E}xt^i_Y(\pi_*\OO_X,\omega_Y)$ 是\trCoh{}的（\mbox{练习~29.2.F}），故只须证明：当 $m\gg 0$ 时，$\mathscr{E}xt^i_Y(\pi_*\OO_X,\omega_Y)\otimes\OO(m)$ 没有非零的整体截面（因为对 $Y$ 上任意\trCoherentSheaf{} $\mathscr{G}$，当 $m\gg 0$ 时 $\mathscr{G}(m)$ 由整体截面生成——见\trSerre{}定理~A，\mbox{定理~16.1.5}）。由\mbox{练习~29.2.J}，
\[
\mathrm{H}^j(Y,\mathscr{E}xt^i_Y(\pi_*\OO_X,\omega_Y)(m))=\mathrm{H}^j(Y,\mathscr{E}xt^i_Y(\pi_*\OO_X,\omega_Y(m))).
\]
若 $j>0$，则当 $m\gg 0$ 时，由\trSerre{}消失，$\mathrm{H}^j(Y,\mathscr{E}xt^i_Y(\pi_*\OO_X,\omega_Y)(m))=0$。于是，由对 $\mathrm{Ext}$ 的局部到整体\trSpectralSequence（\mbox{练习~29.2.M}），
\[
\mathrm{H}^0(Y,\mathscr{E}xt^i_Y(\pi_*\OO_X,\omega_Y(m)))=\mathrm{Ext}^i_Y(\pi_*\OO_X,\omega_Y(m)).
\]
再用一次\mbox{练习~29.2.J}，继而由 $Y$ 上对 $\mathrm{Ext}$ 的函子性\trSerreDuality，
\[
\mathrm{Ext}^i_Y(\pi_*\OO_X,\omega_Y(m))=\mathrm{Ext}^i_Y(\pi_*\OO_X(-m),\omega_Y)=\mathrm{H}^{N-i}(Y,\pi_*\OO_X(-m)),
\]
它在 $N-i>n$ 时为 $0$，因为\trProjective\trScheme{}上一个\trQCSheaf{}的\trCohomology{}在高于该层\trSupport{}之\trDimension{}的次数上消失（维数\trCohomology{}消失，\mbox{定理~18.2.6}）。\footnote{译注：原书上述第二、三式中的「$\mathrm{Ext}^i_Y(\pi_*\OO_X,\omega_Y(m))$」均多印一枚右括号（作「$\mathrm{Ext}^i_Y(\pi_*\OO_X,\omega_Y(m)))$」），当为排印之误，径改。——译者注}\hfill$\square$

\mbox{定理~29.4.3} 的证明中最困难的一步，是下面的引理。

\textbf{29.4.5. \trhLemma。}\ ——\ \emph{设 $\pi\colon X\to Y$ 是余维数为 $r$ 的纯维\trProjective{} $k$-\trScheme{}之间的一个\trClosedEmbedding，且 $Y$ 满足对 $\mathrm{Ext}$ 的函子性\trSerreDuality。那么，我们有随 $\mathscr{F}\in\mathit{Coh}_X$ 函子的同构}
\begin{equation*}
\tag{29.4.5.1}\label{eq:hom-ext-iso-lemma}
\mathrm{Hom}_X(\mathscr{F},\pi^{-1}\mathscr{E}xt^r_Y(\pi_*\OO_X,\omega_Y))\overset{\sim}{\longleftrightarrow}\mathrm{Ext}^r_Y(\pi_*\mathscr{F},\omega_Y).
\end{equation*}

\emph{证明。}\quad 选取 $\omega_Y$ 的一个内射分解
\[
0\longrightarrow\omega_Y\longrightarrow\mathscr{I}^0\longrightarrow\mathscr{I}^1\longrightarrow\mathscr{I}^2\longrightarrow\cdots
\]
为了计算 \eqref{eq:hom-ext-iso-lemma} 的右边，我们从这一分解中去掉 $\omega_Y$，然后施加 $\mathrm{Hom}_Y(\pi_*\mathscr{F},\cdot)$，便得到
\[
0\longrightarrow\mathrm{Hom}_Y(\pi_*\mathscr{F},\mathscr{I}^0)\longrightarrow\mathrm{Hom}_Y(\pi_*\mathscr{F},\mathscr{I}^1)\longrightarrow\mathrm{Hom}_Y(\pi_*\mathscr{F},\mathscr{I}^2)\longrightarrow\cdots,
\]
或者（由 $\pi$ 为\trClosedEmbedding{}时 $\pi_*$ 与 $\pi^!$ 在 $\OO$-\trModule{}上的\trAdjointness，\mbox{练习~29.4.B(a)}）：
\[
0\longrightarrow\mathrm{Hom}_X(\mathscr{F},\pi^!\mathscr{I}^0)\longrightarrow\mathrm{Hom}_X(\mathscr{F},\pi^!\mathscr{I}^1)\longrightarrow\mathrm{Hom}_X(\mathscr{F},\pi^!\mathscr{I}^2)\longrightarrow\cdots.
\]
受此启发，考虑复形
\begin{equation*}
\tag{29.4.5.2}\label{eq:pi-bang-complex}
0\longrightarrow\pi^!\mathscr{I}^0\longrightarrow\pi^!\mathscr{I}^1\longrightarrow\pi^!\mathscr{I}^2\longrightarrow\cdots.
\end{equation*}
首先请注意，\eqref{eq:pi-bang-complex} 确实是一个复形；其次请注意，诸 $\pi^!\mathscr{I}^i$ 都是内射的 $\OO_X$-\trModule（由\mbox{练习~29.4.B(b)}）。

\textbf{29.4.D.（练习）}\quad 证明：\eqref{eq:pi-bang-complex} 在第 $i$ 步处的\trCohomology{}是 $\pi^{-1}\mathscr{E}xt^i_Y(\pi_*\OO_X,\omega_Y)$。

于是，由\mbox{命题~29.4.4}，复形 \eqref{eq:pi-bang-complex} 在第 $r$ 步之前是正合的。

\textbf{29.4.E.（练习）}\quad 证明：存在一个直和分解 $\pi^!\mathscr{I}^r=\mathscr{J}\oplus\mathscr{K}$，使得
\[
0\longrightarrow\pi^!\mathscr{I}^0\longrightarrow\pi^!\mathscr{I}^1\longrightarrow\cdots\longrightarrow\pi^!\mathscr{I}^{r-1}\longrightarrow\mathscr{J}\longrightarrow 0
\]
正合。提示：先做出 $r=1$ 的情形。再一条提示：注意，若
\[
0\longrightarrow\mathscr{I}^0\overset{\alpha}{\longrightarrow}\mathscr{F}
\]
正合，而 $\mathscr{I}^0$ 是内射的，则存在一个「分裂」$\alpha$ 的映射 $\beta\colon\mathscr{F}\to\mathscr{I}^0$，它使我们能把 $\mathscr{F}$ 写成直和 $\mathscr{I}^0\oplus\mathscr{J}^1$。若 $\mathscr{F}$ 进而还是内射的，则 $\mathscr{J}^1$ 也是内射的。这便是一个归纳法的开端。

\textbf{29.4.F.（练习）}\quad （我们沿用上一道练习的记号。）请把 $\ker(\mathscr{K}\to\pi^!\mathscr{I}^{r+1})$ 识别为 $\pi^{-1}\mathscr{E}xt^r_Y(\pi_*\OO_X,\omega_Y)$。

\textbf{29.4.G.（练习）}\quad 把上面的各块拼合起来，以完成\mbox{引理~29.4.5} 的证明。\hfill$\square$

我们现在已经为证明\mbox{定理~29.4.3} 做好了准备。

\textbf{29.4.6.}\ \emph{\mbox{定理~29.4.3} 的证明。}\quad 设 $\mathscr{F}$ 是 $X$ 上的一个\trCoherentSheaf。我们希望描述一个同构
\[
\mathrm{Hom}_X(\mathscr{F},\pi^{-1}\mathscr{E}xt^r_Y(\pi_*\OO_X,\omega_Y))\overset{\sim}{\longleftrightarrow}\mathrm{H}^n(X,\mathscr{F})^{\vee},
\]
它随 $\mathscr{F}$ 函子。由\mbox{引理~29.4.5}，我们有一个函子性的同构
\[
\mathrm{Hom}_X(\mathscr{F},\pi^{-1}\mathscr{E}xt^r_Y(\pi_*\OO_X,\omega_Y))\overset{\sim}{\longleftrightarrow}\mathrm{Ext}^r_Y(\pi_*\mathscr{F},\omega_Y).
\]
由 $Y$ 上对 $\mathrm{Ext}$ 的函子性\trSerreDuality，我们有一个函子性的同构
\[
\mathrm{Ext}^r_Y(\pi_*\mathscr{F},\omega_Y)\overset{\sim}{\longleftrightarrow}\mathrm{H}^{N-r}(Y,\pi_*\mathscr{F})^{\vee}.
\]
由于 $\mathscr{F}$ 是 $X$ 上的\trCoherentSheaf，而 $N-r=n$，右边正是 $\mathrm{H}^n(X,\mathscr{F})^{\vee}$。\hfill$\square$

\textbf{29.4.7. \mbox{定理~29.4.3} 的应用。}\quad 我们首先把\mbox{定理~29.4.3} 应用于 $X$ 是 $Y$ 上的一个\trEffectiveCartierDivisor{}这一特殊情形。我们可以这样计算\trDualizingSheaf{} $\pi^{-1}\mathscr{E}xt^1_Y(\pi_*\OO_X,\omega_Y)$：用 $\pi_*\OO_X$ 的任意一个局部自由分解（在 $Y$ 上）来计算 $\mathscr{E}xt^1_Y(\pi_*\OO_X,\omega_Y)$（\mbox{练习~29.2.E}）。但 $\pi_*\OO_X$ 有一个特别简单的分解，即 $X$ 的\trClosedSubscheme{}正合列~\mbox{(9.1.2.1)}：
\begin{equation*}
\tag{29.4.7.1}\label{eq:cartier-closed-subscheme-seq}
0\longrightarrow\OO_Y(-X)\longrightarrow\OO_Y\longrightarrow\pi_*\OO_X\longrightarrow 0.
\end{equation*}
我们这样计算 $\mathscr{E}xt^{\bullet}(\pi_*\OO_X,\omega_Y)$：截断（去掉 $\pi_*\OO_X$），然后施加 $\mathscr{H}om(\cdot,\omega_Y)$。$\mathscr{E}xt^{\bullet}(\pi_*\OO_X,\omega_Y)$ 是如下复形的上同调
\[
0\longrightarrow\mathscr{H}om_Y(\OO_Y,\omega_Y)\longrightarrow\mathscr{H}om(\OO_Y(-X),\omega_Y)\longrightarrow 0,
\]
也就是说，是
\[
0\longrightarrow\omega_Y\longrightarrow\omega_Y\otimes\OO_Y(X)\longrightarrow 0.
\]
的上同调。

我们立即看到：当 $i\ne 0$、$1$ 时，$\mathscr{E}xt^i(\pi_*\OO_X,\omega_Y)=0$。此外，由\mbox{命题~29.4.4}，$\mathscr{E}xt^0(\pi_*\OO_X,\omega_Y)=0$。（一条不重要的题外话：你可以用它来证明 $\omega_Y$ 没有\trEmbeddedPoint。）

我们现在考虑 $\omega_X=\coker(\omega_Y\to\omega_Y\otimes\OO_Y(X))$。用\trInvertibleSheaf{} $\OO_Y(X)$ 张量 \eqref{eq:cartier-closed-subscheme-seq}，随后再与 $\omega_Y$ 作张量，便得到
\[
\omega_Y\longrightarrow\omega_Y\otimes\OO_Y(X)\longrightarrow\omega_Y\otimes\pi_*(\OO_Y(X)|_X)\longrightarrow 0.
\]
右边的项 $\omega_Y\otimes\pi_*(\OO_Y(X)|_X)$ 常常（多少不太正式地）被写作 $\omega_Y(X)|_X$。于是
\[
\omega_X=\coker(\omega_Y\to\omega_Y\otimes\OO_Y(X))=\omega_Y(X)|_X,
\]
而且这一识别是\emph{\trCanonical{}的}。

我们便证明了下述结果。

\textbf{29.4.8. \trhProposition（\trAdjunctionFormula）。}\ ——\ \emph{设 $Y$ 是\trPureDimension{}为 $n$ 的 \trCM{} \trProjectiveScheme（由\mbox{推论~29.3.14}，它满足对 $\mathrm{Ext}$ 的函子性\trSerreDuality），而 $X$ 是 $Y$ 上的一个\trEffectiveCartierDivisor。（于是由 \S26.2.4，$X$ 是\trCM{}的，从而由\mbox{推论~29.3.14}，它满足对 $\mathrm{Ext}$ 的函子性\trSerreDuality。）则 $\omega_X=\omega_Y(X)|_X$。特别地，若 $\omega_Y$ 是 $Y$ 上的一个\trInvertibleSheaf，则 $\omega_X$ 是 $X$ 上的一个\trInvertibleSheaf。}

作为一个直接的应用，我们有下面的结果。

\textbf{29.4.H.（练习）}\quad 设 $X$ 是 $\mathbb{P}^n$ 中的一个\trCompleteIntersection，由次数为 $d_1$，…，$d_r$ 的诸\trHypersurface{}给出。（注意：由\mbox{命题~26.2.6}，$X$ 是\trCM{}的，从而由\mbox{推论~29.3.14}，它满足对 $\mathrm{Ext}$ 的函子性\trSerreDuality。）证明：$\omega_X\cong\OO_X(-n-1+\sum d_i)$。进而，若 $X$ 还是\trSmooth{}的，证明：$\omega_X\cong\det\Omega_X$。（最后一句话的提示：使用 $\mathscr{K}$ 的\trAdjunctionFormula，即\mbox{练习~21.5.B}。）

但我们还能说得更多。

\textbf{29.4.I.（练习）}\quad 设 $\pi\colon X\hookrightarrow Y$ 是余维数为 $r$ 的一个\trRegularEmbedding，其正规层为 $\mathscr{N}_{X/Y}$（由\mbox{命题~21.2.16(b)}，它是\trLocallyFree{}的）。设 $\mathscr{L}$ 是 $Y$ 上的一个\trInvertibleSheaf。

(a) 证明：当 $i\ne r$ 时，$\mathscr{E}xt^i_Y(\pi_*\OO_X,\mathscr{L})=0$。

(b) 请描述一个\trCanonical{}同构
\[
\pi^{-1}\mathscr{E}xt^r_Y(\pi_*\OO_X,\mathscr{L})\overset{\sim}{\longleftrightarrow}(\det\mathscr{N}_{X/Y})\otimes_{\OO_X}\mathscr{L}|_X.
\]

两部分共同的提示：先处理 $\mathscr{L}=\OO$ 的情形，然后注意到：假如你已经细心而正确地解出了 (b)，这个问题便是局部的。

（注意：由于\mbox{练习~29.4.I} 与对偶性并无显式关联，我们对 $Y$ 不作任何\trProjective{}性假设；这完全是一个局部的问题。）由\mbox{练习~29.4.I}，我们推出下面的结论。

\textbf{29.4.J.（重要练习）}\quad 设 $X$ 是 $\mathbb{P}^n_k$ 中余维数为 $r$ 的一个\trRegularEmbedding。（此时 $X$ 由\mbox{命题~26.2.6} 是\trCM{}的，从而由\mbox{推论~29.3.14} 满足对 $\mathrm{Ext}$ 的函子性\trSerreDuality。）证明：
\[
\omega_X=\pi^{-1}\mathscr{E}xt^r_{\mathbb{P}^n}(\pi_*\OO_X,\omega_{\mathbb{P}^n})\cong(\det\mathscr{N}_{X/\mathbb{P}^n})\otimes\omega_{\mathbb{P}^n}|_X.
\]
特别地，$\omega_X$ 是一个\trInvertibleSheaf。

\textbf{29.4.9.}\ \emph{题外话：Gorenstein 奇点。}\quad 若纯维的 \trCM{} \trProjective{} $k$-\trScheme{} $X$ 的 $\omega_{X/k}$ 是一个\trInvertibleSheaf，我们就说 $X$ 是 \emph{Gorenstein} 的。借助\mbox{命题~29.4.8}，你可以证明：这是一个可以在纯维\trProjective{} $k$-\trScheme{}的各个\trClosedPoint{}处按茎局部检验的条件，正如 \trCM{}性的定义那样。你能写出这样一个定义吗？在正确的普遍性下，Gorenstein 奇点的定义当然并不要求所论的\trScheme{}是射影的。这段题外话将引导你借助一个「\trSlicing{}判别法」，得到一个正确却谈不上优雅的定义。一切「局部\trCompleteIntersection」的 $k$-\trScheme{}都是 Gorenstein 的，而一切 Gorenstein 的 $k$-\trScheme{}按定义都是\trCM{}的。关于 Gorenstein 环与 Gorenstein 概形的一份极佳入门材料，见 [BruH, §3]。

\textbf{29.4.K.（重要练习）}\quad 设 $X$ 是 $\mathbb{P}^n_k$ 的一个\trSmooth{}且纯余维数为 $r$ 的\trSubvariety（因而由\mbox{练习~13.2.M(b)}，它是一个\trRegularEmbedding）。证明：$\omega_X\cong\mathscr{K}_X$。提示：两边都满足\trAdjunctionFormula（关于 $\Omega$ 的\trAdjunctionFormula，见\mbox{练习~21.5.B}）：它们同构于 $(\det\mathscr{N}_{X/Y})\otimes(\omega_{\mathbb{P}^n}|_X)\cong(\det\mathscr{N}_{X/Y})\otimes(\mathscr{K}_{\mathbb{P}^n}|_X)$。\footnote{译注：原书这两处下标作「$X/Y$」，其中 $Y$ 未在本练习中出现（本练习中 $X\subset\mathbb{P}^n_k$），疑当作「$X/\mathbb{P}^n$」，兹照录原书。——译者注}

\textbf{29.4.10.}\ \emph{\trhRemark{}。}\quad 正如久已许诺的那样：把\mbox{推论~29.3.10} 与\mbox{练习~29.4.K} 结合起来，\mbox{定理~18.5.1} 与\mbox{进一步的奇迹~18.5.2} 中给出的那个版本的\trSerreDuality{}便随之得证。


\epigraph{在那些乐于为理解而忧心、追逐完美的人眼中，一部作品永远不会完成——「完成」一词于他们毫无意义——而只是被放弃。而这种放弃——正是它把作品交付给烈火或公众（或许是出于倦怠，或许是出于交稿期限）——在他们看来不过是某种意外罢了，恰似叩门声骤然响起，惊破了他们的梦境……}{——保罗·瓦莱里，[Coh, pp.~7--8]（Tadashi Tokieda 英译）}

\epigraph{好了，我已经给你们讲了两年课，现在就要收场了。在某些方面我想道个歉，在另一些方面则不想。我希望——事实上，我知道——你们当中能有二三十人一直满怀兴奋地听懂了全部内容，并从中获得了乐趣。但我也明白：「讲授的效力微乎其微，除非身处那种讲授几乎纯属多余的幸福境遇。」所以，对那二三十位听懂了全部内容的人，请允许我说一句：我其实什么也没做，不过是把那些东西展示给了你们。至于其余的人，倘若我让你们恨上了这门学科，我很抱歉。我以前从未教过基础物理，在此致歉。我只希望自己没有给你们造成严重的困扰，也希望你们不要就此离开这件令人兴奋的事业。我希望将来能有别人以一种不让你们消化不良的方式把这门课教给你们，而有朝一日你们终究会发现：它并不像看上去那么可怕。}{——R.~费曼，[FLS, Epilogue]}
